% Part: turing-machines
% Chapter: undecidability
% Section: representing-tms

\documentclass[../../../include/open-logic-section]{subfiles}

\begin{document}

\olfileid{tur}{und}{rep}
\olsection{تمثيل آلات تورنغ}

\begin{explain}
تمثيل آلة تورنغ وسلوكها بـ!!a{sentence} من الرتبة الأولى يقتضي
اختيار لغة تلائم هذا الغرض. ولها صنفان من التعبيرات:
!!{predicate}s تصف تهيئات الآلة، وتعبيرات نرقم بها خطوات التنفيذ،
أي «اللحظات»، ومواضع الشريط.

فنضع صنفين من ال!!{predicate}s الثنائية المحل: لكل حالة~${\OLId{latin-ordinary}{q}}$
!!a{predicate}~$\Obj Q_{\OLId{latin-ordinary}{q}}$، ولكل رمز شريط~$\sigma$
!!a{predicate}~$\Obj S_\sigma$. فبالصنف الأول نبيّن حالة~${\OLId{latin-looped}{M}}$
وموضع رأسها، وبالثاني نبيّن ما في الشريط من رموز.

ولا بد، لتعيين موضع الرأس وعدد الخطوات المنفذة، من تعبير عن
الأعداد. فنتخذ لذلك !!a{constant}~$\Obj 0$ ودالة خلف أحادية
المحل~$\prime$. واصطلاحنا أن نكتب علامة هذه الدالة \emph{بعد}
وسيطها، مع إسقاط القوسين.

ويقابل كل عدد~${\OLId{latin-ordinary}{n}}$ حد مخصوص~$\num{{\OLId{latin-ordinary}{n}}}$، نسميه
\emph{العدد القياسي} لـ~${\OLId{latin-ordinary}{n}}$، وبه نمثل العدد في~$\Lang L_{\OLId{latin-looped}{M}}$.
فـ$\num{{\OLMathNumeral{0}}}$ هو $\Obj 0$، و$\num{{\OLMathNumeral{1}}}$ هو $\Obj 0'$،
و$\num{{\OLMathNumeral{2}}}$ هو $\Obj 0''$، وعلى هذا القياس. وحد ذلك بدقة:
\begin{align*}
\num{{\OLMathNumeral{0}}} & = \Obj 0 \\
\num{{\OLId{latin-ordinary}{n}}+{\OLMathNumeral{1}}} &= \num{{\OLId{latin-ordinary}{n}}}'
\end{align*}
فالحد $\num{{\OLMathNumeral{0}}}$، وهو $\Obj 0$، اسم للخانة القصوى في اليسار،
واسم أيضًا للزمن السابق لأول خطوة تنفيذ، أي زمن التهيئة الابتدائية.
والحد $\num{{\OLMathNumeral{1}}}$، وهو $\Obj 0'$، اسم للخانة التالية لها يمينًا،
وللزمن الذي يلي أول خطوة تنفيذ؛ وهكذا فيما بعد.

ونزيد !!a{predicate}~$<$، نستعمله لترتيب المواضع فيكون معناه
«إلى يسار»، ولترتيب خطوات التنفيذ فيكون معناه «قبل».

ثم نورد، بعد تعيين اللغة، «بديهيات» $!T({\OLId{latin-looped}{M}},{\OLId{latin-ordinary}{w}})$؛ وهي
ال!!{sentence}s التي يدل مجموعها على سلوك~${\OLId{latin-looped}{M}}$ عند الدخل~${\OLId{latin-ordinary}{w}}$.
فمنها !!{sentence}s تفرض شروطًا على $\Obj 0$ و$\prime$ و$<$،
ومنها !!{sentence}s لتهيئة الدخل، ومنها !!{sentence}s تبين
تهيئة~${\OLId{latin-looped}{M}}$ الحاصلة عقب تنفيذ تعليمة بعينها.
\end{explain}

\begin{defn}
  \ollabel{defn:tm-descr}
لتكن ${\OLId{latin-looped}{M}}=\tuple{{\OLId{latin-looped}{Q}},\Sigma,{\OLId{latin-ordinary}{q}}_{\OLMathNumeral{0}},\delta}$ آلة تورنغ. نعيّن اللغة
$\Lang L_{\OLId{latin-looped}{M}}$ بالأجزاء الآتية:
\begin{enumerate}
\item !!{predicate} ثنائي المحل $ \Obj Q_{\OLId{latin-ordinary}{q}}({\OLId{latin-ordinary}{x}},{\OLId{latin-ordinary}{y}})$ لكل حالة~${\OLId{latin-ordinary}{q}}\in {\OLId{latin-looped}{Q}}$.
  والمقصود بـ$\Obj Q_{\OLId{latin-ordinary}{q}}(\num{{\OLId{latin-ordinary}{m}}},\num{{\OLId{latin-ordinary}{n}}})$ أن ${\OLId{latin-looped}{M}}$ تكون، بعد ${\OLId{latin-ordinary}{n}}$
  خطوة، في الحالة~${\OLId{latin-ordinary}{q}}$ ورأسها على الخانة رقم~${\OLId{latin-ordinary}{m}}$.
\item !!{predicate} ثنائي المحل $\Obj S_\sigma({\OLId{latin-ordinary}{x}},{\OLId{latin-ordinary}{y}})$ لكل رمز
  $\sigma\in\Sigma$. والمقصود بـ$\Obj S_\sigma(\num{{\OLId{latin-ordinary}{m}}},\num{{\OLId{latin-ordinary}{n}}})$
  أن الخانة رقم~${\OLId{latin-ordinary}{m}}$ تحتوي، بعد ${\OLId{latin-ordinary}{n}}$ خطوة، الرمز~$\sigma$.
\item !!a{constant}~$\Obj 0$.
\item !!a{function} أحادية المحل~$\prime$.
\item !!a{predicate} ثنائي المحل~$<$.
\end{enumerate}
\end{defn}

وهذه ال!!{sentence}s الواصفة لعمل الآلة~${\OLId{latin-looped}{M}}$ عند الدخل
${\OLId{latin-ordinary}{w}}=\sigma_{{\OLId{latin-ordinary}{i}}_{\OLMathNumeral{1}}}\dots\sigma_{{\OLId{latin-ordinary}{i}}_{\OLId{latin-ordinary}{k}}}$:
\begin{enumerate}
\item بديهيات الأعداد والعلاقة~$<$:
\begin{enumerate}
\item !!^a{sentence} مفادها أن كل عدد دون خلفه:
\[
\lforall[{\OLId{latin-ordinary}{x}}][{\OLId{latin-ordinary}{x}} < {\OLId{latin-ordinary}{x}}']
\]
\item !!^a{sentence} تنص على تعدي $<$:
\[
\lforall[{\OLId{latin-ordinary}{x}}][\lforall[{\OLId{latin-ordinary}{y}}][\lforall[{\OLId{latin-ordinary}{z}}][
      (({\OLId{latin-ordinary}{x}} < {\OLId{latin-ordinary}{y}} \land {\OLId{latin-ordinary}{y}} < {\OLId{latin-ordinary}{z}}) \lif {\OLId{latin-ordinary}{x}} < {\OLId{latin-ordinary}{z}})]]]
\]
\end{enumerate}
\item بديهيات تهيئة الدخل:
\begin{enumerate}
\item عند تمام ${\OLMathNumeral{0}}$ من الخطوات، أي قبل التنفيذ، تكون ${\OLId{latin-looped}{M}}$ على الحالة
  الابتدائية~${\OLId{latin-ordinary}{q}}_{\OLMathNumeral{0}}$ ورأسها ماسح للخانة رقم~${\OLMathNumeral{1}}$:
\[
\Obj Q_{{\OLId{latin-ordinary}{q}}_{\OLMathNumeral{0}}}(\num{{\OLMathNumeral{1}}},\num{{\OLMathNumeral{0}}})
\]
\item الخانات ${\OLId{latin-ordinary}{k}}+{\OLMathNumeral{1}}$ الأولى تحمل الرموز $\TMendtape$،
  و$\sigma_{{\OLId{latin-ordinary}{i}}_{\OLMathNumeral{1}}}$،~\dots، و$\sigma_{{\OLId{latin-ordinary}{i}}_{\OLId{latin-ordinary}{k}}}$:
\[
\Obj S_\TMendtape(\num{{\OLMathNumeral{0}}}, \num{{\OLMathNumeral{0}}}) \land
\Obj S_{\sigma_{{\OLId{latin-ordinary}{i}}_{\OLMathNumeral{1}}}}(\num{{\OLMathNumeral{1}}}, \num{{\OLMathNumeral{0}}}) \land
\dots \land
\Obj S_{\sigma_{{\OLId{latin-ordinary}{i}}_{\OLId{latin-ordinary}{k}}}}(\num{{\OLId{latin-ordinary}{k}}}, \num{{\OLMathNumeral{0}}})
\]
\item وسائر الشريط خال من غير علامة الفراغ:
\[
\lforall[{\OLId{latin-ordinary}{x}}][(\num{{\OLId{latin-ordinary}{k}}} < {\OLId{latin-ordinary}{x}} \lif \Obj S_\TMblank({\OLId{latin-ordinary}{x}}, \num{{\OLMathNumeral{0}}}))]
\]
\end{enumerate}
\item بديهيات الانتقال بين التهيئتين المتعاقبتين:

ولنرمز بـ$!A({\OLId{latin-ordinary}{x}},{\OLId{latin-ordinary}{y}})$ إلى اقتران ال!!{sentence}s كلها مما صورته
\[
\lforall[{\OLId{latin-ordinary}{z}}][
  ((({\OLId{latin-ordinary}{z}} < {\OLId{latin-ordinary}{x}} \lor {\OLId{latin-ordinary}{x}} < {\OLId{latin-ordinary}{z}}) \land \Obj S_\sigma({\OLId{latin-ordinary}{z}}, {\OLId{latin-ordinary}{y}}))
  \lif \Obj S_\sigma({\OLId{latin-ordinary}{z}}, {\OLId{latin-ordinary}{y}}'))]
\]
وذلك لكل $\sigma\in\Sigma$. فمعنى $!A(\num{{\OLId{latin-ordinary}{m}}},\num{{\OLId{latin-ordinary}{n}}})$ أن
الشريط بعد ${\OLId{latin-ordinary}{n}}+{\OLMathNumeral{1}}$ خطوة، باستثناء الخانة رقم~${\OLId{latin-ordinary}{m}}$، لم يتغير عما
كان عليه بعد ${\OLId{latin-ordinary}{n}}$ خطوة.
\begin{enumerate}
\item \ollabel{rep-right} نقابل كل تعليمة
$\delta({\OLId{latin-ordinary}{q}}_{\OLId{latin-ordinary}{i}},\sigma)=\tuple{{\OLId{latin-ordinary}{q}}_{\OLId{latin-ordinary}{j}},\sigma',\TMright}$ بال!!{sentence}:
\begin{align*}
& \lforall[{\OLId{latin-ordinary}{x}}][\lforall[{\OLId{latin-ordinary}{y}}][(
   (\Obj Q_{{\OLId{latin-ordinary}{q}}_{\OLId{latin-ordinary}{i}}}({\OLId{latin-ordinary}{x}}, {\OLId{latin-ordinary}{y}}) \land \Obj S_{\sigma}({\OLId{latin-ordinary}{x}}, {\OLId{latin-ordinary}{y}})) \lif {}]] \\
&\qquad   (\Obj Q_{{\OLId{latin-ordinary}{q}}_{\OLId{latin-ordinary}{j}}}({\OLId{latin-ordinary}{x}}', {\OLId{latin-ordinary}{y}}') \land \Obj S_{\sigma'}({\OLId{latin-ordinary}{x}}, {\OLId{latin-ordinary}{y}}') \land
!A({\OLId{latin-ordinary}{x}}, {\OLId{latin-ordinary}{y}})))
\end{align*}
ومفاد ذلك: إذا مضت~${\OLId{latin-ordinary}{y}}$ خطوة وكانت الآلة في الحالة~${\OLId{latin-ordinary}{q}}_{\OLId{latin-ordinary}{i}}$ على
الخانة~${\OLId{latin-ordinary}{x}}$ ورمزها~$\sigma$، فإنها عند تمام ${\OLId{latin-ordinary}{y}}+{\OLMathNumeral{1}}$ خطوة تكون على
الخانة~${\OLId{latin-ordinary}{x}}+{\OLMathNumeral{1}}$ في الحالة~${\OLId{latin-ordinary}{q}}_{\OLId{latin-ordinary}{j}}$، وتحمل الخانة~${\OLId{latin-ordinary}{x}}$ الرمز
$\sigma'$؛ وأما ما عدا~${\OLId{latin-ordinary}{x}}$ من الخانات فيبقى على ما كان فيه
من الرموز بعد~${\OLId{latin-ordinary}{y}}$ خطوة.

\item \ollabel{rep-left} نقابل كل تعليمة
$\delta({\OLId{latin-ordinary}{q}}_{\OLId{latin-ordinary}{i}},\sigma)=\tuple{{\OLId{latin-ordinary}{q}}_{\OLId{latin-ordinary}{j}},\sigma',\TMleft}$ بال!!{sentence}:
\begin{align*}
& \lforall[{\OLId{latin-ordinary}{x}}][\lforall[{\OLId{latin-ordinary}{y}}][
    ((\Obj Q_{{\OLId{latin-ordinary}{q}}_{\OLId{latin-ordinary}{i}}}({\OLId{latin-ordinary}{x}}', {\OLId{latin-ordinary}{y}}) \land \Obj S_{\sigma}({\OLId{latin-ordinary}{x}}', {\OLId{latin-ordinary}{y}})) \lif {}]]\\
& \qquad   (\Obj Q_{{\OLId{latin-ordinary}{q}}_{\OLId{latin-ordinary}{j}}}({\OLId{latin-ordinary}{x}}, {\OLId{latin-ordinary}{y}}') \land \Obj S_{\sigma'}({\OLId{latin-ordinary}{x}}', {\OLId{latin-ordinary}{y}}') \land
!A({\OLId{latin-ordinary}{x}}', {\OLId{latin-ordinary}{y}}))) \land {}\\
& \lforall[{\OLId{latin-ordinary}{y}}][((\Obj Q_{{\OLId{latin-ordinary}{q}}_{\OLId{latin-ordinary}{i}}}(\num{{\OLMathNumeral{0}}}, {\OLId{latin-ordinary}{y}}) \land \Obj S_{\sigma}(\num{{\OLMathNumeral{0}}},
    {\OLId{latin-ordinary}{y}})) \lif {}]\\
& \qquad (\Obj Q_{{\OLId{latin-ordinary}{q}}_{\OLId{latin-ordinary}{j}}}(\num{{\OLMathNumeral{0}}}, {\OLId{latin-ordinary}{y}}') \land \Obj S_{\sigma'}(\num{{\OLMathNumeral{0}}},
  {\OLId{latin-ordinary}{y}}') \land !A(\num{{\OLMathNumeral{0}}}, {\OLId{latin-ordinary}{y}})))
\end{align*}
وتأمل وجه هذه الصياغة: لم نقل «إذا كانت الآلة تمسح الخانة~${\OLId{latin-ordinary}{x}}$،
~\dots»، بل «إذا كانت تمسح الخانة~${\OLId{latin-ordinary}{x}}+{\OLMathNumeral{1}}$، \dots». فإنها بالتحرك
يسارًا تصير في الخطوة التالية إلى الخانة~${\OLId{latin-ordinary}{x}}$، في حين تقع الكتابة
على الخانة~${\OLId{latin-ordinary}{x}}+{\OLMathNumeral{1}}$. وإنما عدلنا إلى هذا التعبير لأن اللغة لا تتضمن
طرحًا ولا دالة سلف.

ولما كانت الأعداد ذات الصورة ${\OLId{latin-ordinary}{x}}+{\OLMathNumeral{1}}$ هي ${\OLMathNumeral{1}}$، و${\OLMathNumeral{2}}$،~\dots،
بقيت حالة الخانة~${\OLMathNumeral{0}}$ خارج هذا الشرط: فقد يؤمر الرأس فيها بالتحرك
يسارًا، ولا موضع له إلى اليسار، فيبقى مكانه. ولهذه الحالة
وضعنا المقترن الثاني. فهو يقول: إذا كانت الآلة بعد ${\OLId{latin-ordinary}{y}}$ خطوة
تمسح الخانة~${\OLMathNumeral{0}}$ في الحالة~${\OLId{latin-ordinary}{q}}_{\OLId{latin-ordinary}{i}}$، وكان في الخانة~${\OLMathNumeral{0}}$
الرمز~$\sigma$، فإنها بعد ${\OLId{latin-ordinary}{y}}+{\OLMathNumeral{1}}$ خطوة تظل على الخانة~${\OLMathNumeral{0}}$،
وتصير إلى الحالة~${\OLId{latin-ordinary}{q}}_{\OLId{latin-ordinary}{j}}$، ويكون في الخانة~${\OLMathNumeral{0}}$ الرمز
$\sigma'$؛ وأما الخانات غير~${\OLMathNumeral{0}}$ فتبقى رموزها كما كانت بعد~${\OLId{latin-ordinary}{y}}$
خطوة.
\item \ollabel{rep-stay} نقابل كل تعليمة
$\delta({\OLId{latin-ordinary}{q}}_{\OLId{latin-ordinary}{i}},\sigma)=\tuple{{\OLId{latin-ordinary}{q}}_{\OLId{latin-ordinary}{j}},\sigma',\TMstay}$ بال!!{sentence}:
\begin{align*}
& \lforall[{\OLId{latin-ordinary}{x}}][\lforall[{\OLId{latin-ordinary}{y}}][(
   (\Obj Q_{{\OLId{latin-ordinary}{q}}_{\OLId{latin-ordinary}{i}}}({\OLId{latin-ordinary}{x}}, {\OLId{latin-ordinary}{y}}) \land \Obj S_{\sigma}({\OLId{latin-ordinary}{x}}, {\OLId{latin-ordinary}{y}})) \lif {}]] \\
&\qquad   (\Obj Q_{{\OLId{latin-ordinary}{q}}_{\OLId{latin-ordinary}{j}}}({\OLId{latin-ordinary}{x}}, {\OLId{latin-ordinary}{y}}') \land \Obj S_{\sigma'}({\OLId{latin-ordinary}{x}}, {\OLId{latin-ordinary}{y}}') \land
!A({\OLId{latin-ordinary}{x}}, {\OLId{latin-ordinary}{y}})))
\end{align*}
\end{enumerate}
\end{enumerate}
ونسمي $!T({\OLId{latin-looped}{M}},{\OLId{latin-ordinary}{w}})$ اقتران ال!!{sentence}s المتقدمة جميعًا، الخاصة
بالآلة~${\OLId{latin-looped}{M}}$ والدخل~${\OLId{latin-ordinary}{w}}$.

وأما توقف~${\OLId{latin-looped}{M}}$ في نهاية المطاف، فنعبر عنه ب!!{sentence}
مفادها أن دالة الانتقال تصير غير معرّفة بعد عدد من الخطوات.
فلتكن ${\OLId{latin-looped}{X}}$ مجموعة الأزواج $\tuple{{\OLId{latin-ordinary}{q}},\sigma}$ التي تكون فيها
$\delta({\OLId{latin-ordinary}{q}},\sigma)$ غير معرّفة؛ ولتكن $!E({\OLId{latin-looped}{M}},{\OLId{latin-ordinary}{w}})$ ال!!{sentence}
\[
\lexists[{\OLId{latin-ordinary}{x}}][\lexists[{\OLId{latin-ordinary}{y}}][(\bigvee_{\tuple{{\OLId{latin-ordinary}{q}}, \sigma} \in
      {\OLId{latin-looped}{X}}}(\Obj Q_{\OLId{latin-ordinary}{q}}({\OLId{latin-ordinary}{x}}, {\OLId{latin-ordinary}{y}}) \land \Obj S_\sigma({\OLId{latin-ordinary}{x}}, {\OLId{latin-ordinary}{y}})))]]
\]

وإن خصصنا للآلة حالة توقف~${\OLId{latin-ordinary}{h}}$، اختصر التعبير؛ فتكون
ال!!{sentence}~$!E({\OLId{latin-looped}{M}},{\OLId{latin-ordinary}{w}})$
\[
\lexists[{\OLId{latin-ordinary}{x}}][\lexists[{\OLId{latin-ordinary}{y}}][\Obj Q_{\OLId{latin-ordinary}{h}}({\OLId{latin-ordinary}{x}}, {\OLId{latin-ordinary}{y}})]]
\]
دالةً على توقف الآلة في نهاية المطاف.

\begin{prop}
\ollabel{prop:mlessk}
من ${\OLId{latin-ordinary}{m}}<{\OLId{latin-ordinary}{k}}$ يلزم $!T({\OLId{latin-looped}{M}},{\OLId{latin-ordinary}{w}})\Entails\num{{\OLId{latin-ordinary}{m}}}<\num{{\OLId{latin-ordinary}{k}}}$.
\end{prop}

\begin{proof}
برهانه متروك للتمرين.
\end{proof}

\begin{prob}
برهن \olref[tur][und][rep]{prop:mlessk}.
(إرشاد: اجعل الاستقراء على ${\OLId{latin-ordinary}{k}}-{\OLId{latin-ordinary}{m}}$.)
\end{prob}

\end{document}
